\section{Lie Groups}

\bd [Lie Group]
A \textbf{Lie group}\index{Lie group} is a group $(G,\bullet)$, where $G$ is a smooth manifold and the maps
\bi{rrCl}
\mu \cl & G\times G & \to & G\\
& (g_1,g_2) & \mapsto & g_1\bullet g_2
\ei
and
\bi{rrCl}
i \cl & G & \to & G\\
& g & \mapsto & g^{-1}
\ei
are both smooth. Note that $G\times G$ inherits a smooth atlas from the smooth atlas of $G$.
\ed

\bd [Dimension Of Lie Group]
The \textbf{dimension} of a Lie group $(G,\bullet)$ is the dimension of $G$ as a manifold.
\ed

\be
\ben[label=\alph*)]
\item Consider $(\R^n,+)$, where $\R^n$ is understood as a smooth $n$-dimensional manifold. This is a commutative (or abelian) Lie group (since $\bullet$ is commutative), often called the $n$-dimensional translation group.

\item Let $S^1:=\{z\in\C\mid |z|=1\}$ and let $\cdot$ be the usual multiplication of complex numbers. Then $(S^1,\cdot)$ is a commutative Lie group usually denoted $\mathrm{U}(1)$.

\item As we discussed in the chapter of vector spaces, $\mathrm{Aut}(V):=\{ \phi\cl V \xrightarrow{\sim} V \mid \det \phi \neq 0 \}$ is the set of all of linear isomorphisms on $V$ that we denoted as $GL(V)$. To make it more concrete, we can take as the vector space $V = \R^n$ hence $\mathrm{GL}(n,\R)=\{\phi\cl\R^n\xrightarrow{\sim}\R^n \mid \det \phi \neq 0 \}$. This set can be endowed with the structure of a smooth $n^2$-dimensional manifold, by noting that there is a bijection between linear maps $\phi\cl\R^n\xrightarrow{\sim}\R^n$ and $\R^{2n}$. Then, $(\mathrm{GL}(n,\R),\circ)$ is a Lie group called the \emph{general linear group}.
\een
\ee

\bd [Lie Group Homomorphism]
Let $(G,\bullet)$ and $(H,\circ)$ be Lie groups. A map $\phi\cl G \to H$ is \textbf{Lie group homomorphism} if it is a group homomorphism and a smooth map.
\ed

\bd [Lie Group Isomorphism]
A \textbf{Lie group isomorphism}\index{isomorphism!of Lie groups} is a Lie group homomorphism which is also a diffeomorphism of the underlying manifolds. 
\ed

\section{The Left Translation Map}

To every element of a Lie group there is associated a special map. Note that everything we will do here can be done equivalently by using right translation maps. 

\bd [Left Translation]
Let $(G,\bullet)$ be a Lie group and let $g\in G$. The map
\bi{rrCl}
\ell_g \cl & G & \to & G\\
& h & \mapsto & \ell_g(h):=g\bullet h \equiv gh
\ei
is called the \textbf{left translation}\index{left translation} by $g$.
\ed

One might thing that this is an overkill of notation since we already had the operation between two element from the group structure. However the left translation map is different, since we first have to fix an element of the group $g$ (hence the index in $\ell_g$) and then apply this element to the whole group (a.k.a to each element of the group).\\ 

If there is no danger of confusion, we usually suppress the $\bullet$ notation.  
\bp
Let $G$ be a Lie group. For any $g\in G$, the left translation map $\ell_g\cl G \to G$ is a isomorphism.
\ep

\bq
Let $h,h'\in G$. Then, we have
\bse
\ell_g(h)=\ell_g(h')\ \Leftrightarrow\ g h = g h' \ \Leftrightarrow\ h=h'.
\ese
Moreover, for any $h\in G$, we have $g^{-1} h\in G$ and
\bse
\ell_g(g^{-1} h) = gg^{-1} h = h.
\ese
Therefore, $\ell_g$ is a bijection on $G$.\\

Note that
\bse
\ell_g = \mu(g,-)
\ese
and since $\mu\cl G\times G \to G$ is smooth by definition, so is $\ell_g$. \\

The inverse map is $(\ell_g)^{-1}=\ell_{g^{-1}}$, since
\bse
\ell_{g^{-1}} \circ \ell_{g} = \ell_{g} \circ \ell_{g^{-1}} = \id_G.
\ese
Then, for the same reason as above with $g$ replaced by $g^{-1}$, the inverse map $(\ell_g)^{-1}$ is also smooth. Hence, the map $\ell_g$ is indeed an isomorphism.
\eq

Note that, $\ell_g$ is not an isomorphism of groups, i.e.\ 
\bse
\ell_g(hh') \neq \ell_g(h)\,\ell_g(h')
\ese
in general. However that does not stop $\ell_g$ from being an isomorphism, which actually means that is a diffeomorphism of the underlying manifolds.\\

Since a lie group is a topological manifold, on top of being a group, this means that at any point $g$ we can define the tangent space, and by following the analysis we did in the previous chapter to define fields on $G$. Recall from the previous chapter than once we have a diffeomorphism $\phi$ between two manifolds $M$ and $N$, we can define the push-forward of a vector field X as
\bse
\phi_* (X) (\phi(p)) := (\phi_*)_p (X_p)
\ese

where $X_p$ is the vector created by the field $X$ at point $p$. \\

Coming in our case, we just showed that the map $\ell_g\cl G\to G$ is a diffeomorphism so we can push-forward any vector field $X$ on $G$ to another vector field (again on $G$ since the maps is between the same manifold). So in our case $\phi_* (X) = (\ell_g)_* (X)$ and for any point $h$ in $G$: $\ell_g (h) = gh$ so the push-forward equation reads 
\bse
{\ell_{g}}_{*} (X) (gh):= ({\ell_{g}}_{*})_{h} (X_h)
\ese

\section{The Lie Algebra Of A Lie Group}

In Lie theory, we are typically not interested in general vector fields, but rather on special class of vector fields which are invariant under the induced push-forward of the left translation maps $\ell_g$.

\bd [Left Invariant Vector Fields]
Let $G$ be a Lie group. A vector field $X\in\Gamma(TG)$ is said to be \textbf{left-invariant}\index{left-invariant vector field} if
\bse
\forall \, g \in G  : \ {\ell_{g}}_{*} (X)= X.
\ese

\vspace{5pt}

Equivalently, we can require this to hold pointwise
\bse
\forall \, g,h \in G : \ ({\ell_{g}}_{*})_{h} (X_h)= X_{gh}.
\ese

where basically what this means is that if we take the vector $X_h$ at $h$ and the vector $X_{gh}$ at $gh$, then by pushing forward $X_h$ with $({\ell_{g}}_{*})_{h}$ (hence bringing it to the point $gh$), the final result should be equal to the vector $X_{gh}$. 
\ed

We can manipulate a bit the pointwise formulation to yield another reformulation. Since both sides are vectors we can let them act on a function $f$ 
\bse
({\ell_{g}}_{*})_{h} (X_h) f = X_{gh} f
\ese

By using the definition of a push-forward of a vector $(\phi_*)_p (X_p) f := X_p(f \circ \phi)$ the left part of the equation reads
\bse
({\ell_{g}}_{*})_{h} (X_h) f = X_h (f \circ \ell_g) = X (h) (f \circ \ell_g)  = X (f \circ \ell_g) (h)
\ese

\vspace{5pt}

The right part can be manipulated as follows
\bse
X_{gh} f = X(gh) f = (Xf) (gh) = (Xf)(\ell_g(h)) =(Xf) \circ \ell_g (h)
\ese

By substituting both final expressions back to the original one and discarding the point $h$ since they must be true for any $h$ we obtain the last reformulation of the push-forward
\bse
X (f \circ \ell_g) = X(f) \circ \ell_g.
\ese

Once again, let's emphasize that this equations holds only for left invariant vector fields and not all vector fields $X$ of $\Gamma(TG)$.

\bd [$\mathcal{L}(G)$ (As A Set)]
We denote the set of all left-invariant vector fields on $G$ as $\mathcal{L}(G)$. 
\ed

Of course,
\bse
\mathcal{L}(G)\se\Gamma(TG)
\ese

\vspace{3pt}

but, in fact, more is true. Recall that we equipped $(\Gamma(TG)$ with two operations and we showed that $(\Gamma(TG),+,\cdot)$ is in fact a $\mathcal{C}^\infty(G)$-submodule. One can check that $\mathcal{L}(G)$ is closed under 
\bi{c}
+\cl \mathcal{L}(G)\times \mathcal{L}(G) \to \mathcal{L}(G)\\
\cdot  \cl \mathcal{C}^\infty(G)\times \mathcal{L}(G) \to \mathcal{L}(G),
\ei
therefore, $\mathcal{L}(G)$ is a $\mathcal{C}^\infty(G)$-submodule of $\Gamma(TG)$. \\

On top of that we said that $(\Gamma(TG),+,\cdot)$ can also be seen as an $\R$-vector space. Up to now, we have refrained from thinking of $\Gamma(TG)$ as an $\R$-vector space since it is infinite-dimensional and, even worse, a basis is in general uncountable. A priori, this could be true for $\mathcal{L}(G)$ as well, but we will see that the situation is, in fact, much nicer as $\mathcal{L}(G)$ will turn out to be a finite-dimensional vector space over $\R$ (as an $\R$-vector subspace of $\Gamma(TG)$). Let's see why, since the reason why it is so it's of crucial importance.

\begin{theorem}
Let $G$ be a Lie group with identity element $e\in G$. Then $\mathcal{L}(G)\cong_\mathrm{vec} T_eG$.
\end{theorem}

\bq
We will construct a linear isomorphism $j\cl T_eG\xrightarrow{\sim}\mathcal{L}(G)$. Define
\bi{rrCl}
j \cl & T_eG& \to & \Gamma(TG)\\
& X_{\gamma,e} & \mapsto & j(X_{\gamma,e}),
\ei
where $j(X_{\gamma,e})$ is define as
\bi{rrCl}
j(X_{\gamma,e}) \cl & G& \to & TG\\
& g & \mapsto & j(X_{\gamma,e}) (g) := ({\ell_{g*}})_e (X_{\gamma,e}) \, \forall g \in G
\ei

\br
Recall that we dropped the curve $\gamma$ in the notation of a vector $X_p$ to save space on the notation. However since we are going to need the definition of a vector, that contains the curve, we will write it again as $X_{\gamma,p}$ for now.
\er

Now we have to prove that this is actually a linear isomorphism, and we will do it in steps.

\ben[label=\roman*)]

\item First we need to check that the definition we provided is consistent. Focusing on the left side of the equation, $j(X_{\gamma,e})$ as a field associates to every point $g$ of the manifold $G$ a tangent vector, hence the first part of the definition $j(X_{\gamma,e})(g)$ is a tangent vector of the tangent space $T_gG$.\\

Now we need to see if the second part of the equation is again a vector on  $T_gG$. Recall that under a push-forward a vector of $T_pM$ goes to a vector of $T_{\phi(p)}N$, hence in our case the vector  $X_{\gamma,e}$ of $T_eG$ goes to the vector $({\ell_{g*}})_e X_{\gamma,e}$ which lies on $T_{l_g(e)}G = T_{ge}G = T_gG$. \\

Hence there is consistency in our definition.

\item Second, we show that for any $X_{\gamma,e}\in T_eG$, $j(X_{\gamma,e})$ is a smooth vector field on $G$. It suffices to check that for any $f\in \mathcal{C}^\infty(G)$, we have $j(X_{\gamma,e})(f)\in \mathcal{C}^\infty(G)$. Indeed
\bi{rCl"s}
j(X_{\gamma,e})(g) & := & ({\ell_{g*}})_e (X_{\gamma,e})& (definition of $j$)\\
j(X_{\gamma,e})(g)(f) & = & ({\ell_{g*}})_e (X_{\gamma,e})(f) & (act on  a function)\\
& = &  X_{\gamma,e}(f\circ\ell_g) & (definition of push-forward) \\
& = &  (f\circ\ell_g\circ\gamma)'(0) & (definition of tangent vector)
\ei
The map
\bi{rrClCl}
\varphi \cl & \R\times G &\to & \R &&\\
&(t,g)&\mapsto & \varphi(t,g) & := &(f\circ\ell_{g}\circ\gamma)(t) \\
&&&& = & f(g\gamma(t))
\ei
is a composition of smooth maps, hence it is smooth. Then
\bse
j(X_{\gamma,e})(g)(f) = (\partial_1\varphi)(0,g)
\ese
depends smoothly on $g$ and thus $j(X_{\gamma,e})(f)\in \mathcal{C}^\infty(G)$.

\item We proved that $j(X_{\gamma,e})$ is indeed a smooth vector field, however now we need to prove that it is a left invariant vector field since it is an element of $\Gamma(TG)$. To prove that we need to shows that it's invariant under left translations. Let $g,h\in G$. Then, for every $X_{\gamma,e}\in T_eG$, we have
\bi{rCl"s}
j(X_{\gamma,e})(g) & := & ({\ell_{g*}})_e (X_{\gamma,e})& (definition of $j$)\\
({\ell_{h*}})_e (j(X_{\gamma,e})(g)) & = & ({\ell_{h*}})_e(({\ell_{g*}})_e(X_{\gamma,e})) & (acting with $({\ell_{h*}})_e$)\\
& = & ({\ell_{gh*}})_e (X_{\gamma,e}) & (one can show that)\\
& = & j(X_{\gamma,e})(gh) & (definition of $j$),
\ei
so $j(X_{\gamma,e})\in \mathcal{L}(G)$. Hence, the map $j$ is really $j\cl T_eG \to \mathcal{L}(G)$.

\item We also need to check the linearity. Let $X_{\gamma,e}, Y_{\gamma,e} \in T_eG$ and $\lambda \in \R$. Then, for any $g\in G$
\bi{rCl"s}
j(\lambda X_{\gamma,e} + Y_{\gamma,e})(g) & = & ({\ell_{g*}})_e (\lambda X_{\gamma,e} + Y_{\gamma,e}) & (definition of $j$)\\
& = & \lambda ({\ell_{h*}})_e( X_{\gamma,e}) +  ({\ell_{h*}})_e (Y_{\gamma,e})\\
& = & \lambda j(X_{\gamma,e})(g) + j(Y_{\gamma,e})(g),
\ei
since the push-forward is an $\R$-linear map. Hence, we have $j\cl T_eG \xrightarrow{\sim} \mathcal{L}(G)$.

\item We also need to check that the map is injective. Let $X_{\gamma,e}, Y_{\gamma,e} \in T_eG$. Then
\bi{rCl"s}
j(X_{\gamma,e}) = j(Y_{\gamma,e}) & \Leftrightarrow & \forall \, g \in G : j(X_{\gamma,e})(g) = j(Y_{\gamma,e})(g)\\
&\Rightarrow & j(X_{\gamma,e})(e) = j(Y_{\gamma,e})(e)\\
 & \Leftrightarrow & ({\ell_{e*}})_e(X_{\gamma,e}) = ({\ell_{e*}})_e(Y_{\gamma,e}) & (definition of $j$)\\
& \Leftrightarrow & X_{\gamma,e}=Y_{\gamma,e} & ($({\ell_{e*}})_e=\id_{TG}$)
\ei
Hence, the map $j$ is injective.

\item Finally we need to check that the map is surjective. Let $X \in \mathcal{L}(G)$. Then, we have
\bi{rCl"s}
j(X_{\gamma,e})(g) & = & ({\ell_{g*}})_e (X_{\gamma,e}) & (definition of $j$)\\
& = & ({\ell_{g*}})_e (X(e)) & ($X(e) = X_{\gamma,e}$)\\
& = & X(ge) & ($X$ is left-invariant) \\
& = & X(g),
\ei
Hence $X=j(X_{\gamma,e})$ and thus $j$ is surjective. 

\een
Therefore, $j\cl T_eG\xrightarrow{\sim}\mathcal{L}(G)$ is indeed a linear isomorphism.
\eq

\bc
It is $\dim \mathcal{L}(G)=\dim G$, hence $\mathcal{L}(G)$ turns out to be a finite-dimensional vector space over $\R$ (as an $\R$-vector subspace of $\Gamma(TG)$).
\ec

So we proved that indeed $\mathcal{L}(G)$ is a finite-dimensional vector space, but as we said there is something more important here. Namely, since $j\cl T_eG\xrightarrow{\sim}\mathcal{L}(G)$ is a linear isomorphism this means that the spaces $\mathcal{L}(G)$ and $T_eG$ are isomorphic which with its turn it means that the map $j$, first of all has an inverse, and most importantly maps each element of $\mathcal{L}(G)$ to exactly one element of $T_eG$. In other words there are as many left invariant fields as there are tangent vectors to the Lie group at the identity, so instead of studying the (quite complicated) left invariant vector fields we can simply study the (quite simpler) tangent vectors at the identity. \\

However, we can go one step further now. Recall from previous chapter that by considering $\Gamma(TM)$ as an infinite-dimensional $R$-vector space, for two vector fields $X,Y\in \Gamma(TM)$, we can define their Lie bracket to be the commutator of the fields
\bse
[X,Y] (f):= X(Y(f))-Y(X(f))
\ese
for any $f\in \mathcal{C}^\infty(M)$. Then $(\Gamma(TG),+,\cdot,[-,-])$ is a Lie algebra. Coming to our case, since $\mathcal{L}(G)$ is a subvector space (and a submodule) of $\Gamma(TM)$ we can inherit the commutator to $\mathcal{L}(G)$ and ask if $\mathcal{L}(G)$ is closed under the commutator so $(\mathcal{L}(G),+,\cdot,[-,-])$ is a subalgebra of $(\Gamma(TG),+,\cdot,[-,-])$. Indeed, this is the case.

\begin{theorem}
Let $G$ be a Lie group. Then $\mathcal{L}(G)$ is a Lie subalgebra of $\Gamma(TG)$.
\end{theorem}
\bq
A Lie subalgebra of a Lie algebra is simply a vector subspace which is closed under the action of the Lie bracket. Therefore, we only need to check that
\bse
\forall \, X,Y \in \mathcal{L}(G) : \ [X,Y]\in \mathcal{L}(G).
\ese
Let $X,Y \in \mathcal{L}(G)$. For any $g\in G$ and $f\in\mathcal{C}^\infty(G)$, we have
\bi{rCl}
[X,Y](f\circ\ell_g) & := & X(Y(f\circ\ell_g))-Y(X(f\circ\ell_g))\\
& = & X(Y(f)\circ\ell_g)-Y(X(f)\circ\ell_g)\\
& = & X(Y(f))\circ\ell_g-Y(X(f))\circ\ell_g\\
& = & \bigl(X(Y(f))-Y(X(f))\bigr)\circ\ell_g\\
& = & [X,Y](f)\circ\ell_g.
\ei
Hence, $[X,Y]$ is left-invariant.
\eq

To summarise, we began with $\mathcal{L}(G)$ as a set of all left invariant vector fields of $G$, which is a subset of $\Gamma(TG)$, then we inherited the $+$ and $\cdot$ of $\Gamma(TG)$ to $\mathcal{L}(G)$ and we showed that it is also a submodule and a subvector space of $\Gamma(TG)$, and finally we inherited the Lie bracket from $\Gamma(TG)$ and we showed that it is also a subalgebra of $\Gamma(TG)$. From now on when we will be referring to $\mathcal{L}(G)$, we will mean its algebra structure.

\bd [The Lie Algebra Of A Lie Group]
Let $G$ be a Lie group. We call the Lie algebra  $\mathcal{L}(G)$ of all left invariant vector fields of $G$ the \textbf{Lie algebra of the Lie group} $G$.
\ed

We already showed that the underlying vector space of the Lie algebra of a Lie group $\mathcal{L}(G)$ is isomorphic to the tangent vector of the Lie group $G$ at the identity $T_eG$. We will now see that the identification of $\mathcal{L}(G)$ and $T_eG$ goes beyond the level of linear isomorphism as vector spaces, as they are isomorphic as algebras. Of course we cannot simply inherit the commutator form $\mathcal{L}(G)$ to $T_eG$ because the commutator is defined for fields and it works because one field can act on another. In the case of $T_eG$ we cannot act with a vector on another vector.  However, we can use the commutator from $\mathcal{L}(G)$ to construct an appropriate Lie bracket on $T_eG$ such that they be isomorphic as algebras.\\

Recall from the Lie algebras chapter in the notes that two algebras are called isomorphic if there exists an isomorphism between them, a.k.a a bijective map $\phi$ such that
\bse
\forall \, x,y\in L_1 : \ \phi([x,y]_{L_1}) = [\phi(x),\phi(y)]_{L_2}.
\ese

Well, we already have a bijective map $j$ so by using the commutator $[-,-]_{\mathcal{L}(G)}$ on $\mathcal{L}(G)$ we can define a Lie bracket for any $X_e, Y_e \in T_eG$ as
\bse
[X_e ,Y_e]_{T_eG} := j^{-1} \bigl( [j(X_e),j(Y_e)]_{\mathcal{L}(G)} \bigr),
\ese

\vspace{5pt}

where $j^{-1} $ works the opposite way, i.e takes a vector field and produces a vector.\\

Equipped with these Lie bracket, by definition, we have
\bse
\mathcal{L}(G)\cong_\mathrm{Lie\, alg}T_eG.
\ese

Hence, given a Lie group we have seen how we can construct its corresponding Lie algebra as the space of left-invariant vector fields and we also showed that this algebra is isomorphic to the algebra of tangent vectors at the identity. We will later explore the opposite direction, i.e.\ given a Lie algebra, we will see how to construct a Lie group whose associated Lie algebra is the one we started from. 

\section{Application: $\SL(2,\C)$ - Part 2}

In the first part of of the application in the previous chapter, we defined the set $\SL(2,\C)$ as a subset of $\C^4:=\C\times\C\times\C\times\C$. Then we showed that:
\bit
\item $\SL(2,\C)$ can be made into a group
\item $\SL(2,\C)$ can be made into a topological space
\item $\SL(2,\C)$ can be made into a topological manifold
\item $\SL(2,\C)$ can be made into a complex differentiable manifold
\eit

Hence we have left with  $\SL(2,\C)$ as a 3-dimensional, complex differentiable manifold.

\subsubsection{$\SL(2,\C)$ As A Lie Group}

We equipped $\SL(2,\C)$ with both a group and a manifold structure. In order to obtain a Lie group structure, we have to check that these two structures are compatible, that is, we have to show that the two maps
\bi{rrCl}
\mu \cl & \SL(2,\C) \times \SL(2,\C) & \to & \SL(2,\C)\\[3pt]
& (\biggl(\begin{matrix} a & b \\ c & d\end{matrix}\biggr) ,\biggl(\begin{matrix} e & f \\ g & h\end{matrix}\biggr))  & \mapsto & \biggl(\begin{matrix} a & b \\ c & d\end{matrix}\biggr) \bullet \biggl(\begin{matrix} e & f \\ g & h\end{matrix}\biggr)
\ei
and 
\bi{rrCl}
i \cl & \SL(2,\C) & \to & \SL(2,\C)\\[3pt]
& \biggl(\begin{matrix} a & b \\ c & d\end{matrix}\biggr)  & \mapsto &\biggl(\begin{matrix} a & b \\ c & d\end{matrix}\biggr)^{-1}
\ei

\vspace{3pt}

are differentiable with respect to the differentiable structure on $\SL(2,\C)$. \\

For instance, for the inverse map $i$, we have to show that the map $y\circ i \circ x^{-1}$ is differentiable in the usual for any pair of charts $(U,x),(V,y)\in \mathscr{A}$. 
\bse
\begin{tikzcd}
U \se\SL(2,\C) \ar[dd,"x"]\ar[rr,"i"]&& V\se \SL(2,\C)\ar[dd,"y"]\\
&&\\
x(U) \se \C^3 \ar[rr,"y\circ i\circ x^{-1}"]&& y(V)\se \C^3
\end{tikzcd}
\ese

\vspace{5pt}

where we remind that
\bi{rrCl}
x^{-1} \cl & x(U) & \to & U\\
& (a,b,c)& \mapsto & \biggl( \begin{matrix} a & b \\ c & \frac{1+bc}{a} \end{matrix}\biggr) .
\ei

and
\bi{rrCl}
y \cl & V & \to & x(V) \se \C\times \C^*\times \C\\
& \biggl( \begin{matrix} a & b \\ c & d \end{matrix}\biggr) & \mapsto & (a,b,d).
\ei

\vspace{3pt}

Since $\SL(2,\C)$ is connected, the differentiability of the transition maps in $\mathscr{A}$ implies that if $y\circ i\circ x^{-1}$ is differentiable for any two given charts, then it is differentiable for all charts in $\mathscr{A}$. Hence, we can simply let $(U,x)$ and $(V,y)$ be the two charts on $\SL(2,\C)$ defined above. Then, we have
\bse
(y\circ i\circ x^{-1}) (a,b,c) = (y\circ i) ( \biggl(\begin{matrix} a & b \\ c & \frac{1+bc}{a}\end{matrix}\biggr) ) = y ( \biggl(\begin{matrix} \frac{1+bc}{a} & -b \\ -c & a\end{matrix}\biggr)) = (\tfrac{1+bc}{a},-b,a)
\ese
which is certainly complex differentiable as a map between open subsets of $\C^3$ (recall that $a\neq 0$ on $x(U)$). We have to do the whole process again for $x\circ i \circ y^{-1}$ and show that is complex differentiable (which it is), hence we conclude that indeed the map $i$ is is complex differentiable. \\

Checking that $\mu$ is complex differentiable is slightly more involved, since we first have to equip $\SL(2,\C) \times \SL(2,\C)$ with a suitable ``product differentiable structure'' and then proceed as above. Once that is done, we can finally conclude that $((\SL(2,\C),\mathcal{O},\mathscr{A}),\bullet)$ is a $3$-dimensional complex Lie group.

\subsubsection{The Lie Algebra $\sl(2,\C)$ Of The Lie Group \texorpdfstring{$\SL(2,\C)$}{SL(2,C)}}

Recall that to every Lie group $G$, there is an associated Lie algebra $\mathcal{L}(G)$ of all left invariant vector fields of $G$ a.k.a
\bse
\mathcal{L}(G) := \{X\in \Gamma(TG)\mid \forall \, g,h\in G : (\ell_g)_*(X|_h)=X|_{gh}\},
\ese

\vspace{5pt}

where the left translation map $\ell_g$ was given by
\bse
\ell_g(h):=g\bullet h \equiv gh
\ese

Coming to our case we have that $G = \SL(2,\C)$ hence the corresponding Lie algebra of $\SL(2,\C)$ usually denoted by small letters $\sl(2,\C)$ is
\bse
\sl(2,\C) := \mathcal{L}(\SL(2,\C)) := \{X\in \Gamma(T\SL(2,\C))\mid \forall \, g,h\in G : (\ell_g)_*(X|_h)=X|_{gh}\},
\ese

where the left translation map at a point $g = \left(\begin{smallmatrix}a & b \\ c & d\end{smallmatrix}\right)$ of $\SL(2,\C)$ is 
\bse
\ell_{\left(\begin{smallmatrix}a & b \\ c & d\end{smallmatrix}\right)} \biggl(\begin{matrix}e & f \\ g  & h\end{matrix}\biggr) = \biggl(\begin{matrix}a & b \\ c & d\end{matrix}\biggr) \biggl(\begin{matrix}e & f \\ g  & h\end{matrix}\biggr) 
\ese

\vspace{5pt}

Now, our goal is to classify $\sl(2,\C)$ as a Lie algebra. The first step we need to do it to find its structure constants. This can be done by considering two vector fields $X,Y\in \Gamma(T\SL(2,\C))$ and computing the commutator
\bse
[X,Y]:= X(Y)-Y(X)
\ese

\vspace{3pt}

However, as we proved earlier, we can always use the fact that the corresponding Lie algebra $\mathcal{L}(G)$ of a Lie group $G$ is isomorphic to the Lie algebra $T_eG$ with Lie bracket
\bse
[X_e ,Y_e]_{T_eG} := j^{-1} \bigl( [j(X_e),j(Y_e)]_{\mathcal{L}(G)} \bigr),
\ese

with
\bse
 j(X_e) (g) := ({\ell_{g*}})_e (X_{e}) \, \forall g \in G
\ese

\vspace{5pt}

Hence, in our case we can study the lie algebra $\sl(2,\C)$ by focusing on the tangent space $T_{\left(\begin{smallmatrix}1 & 0 \\ 0 & 1\end{smallmatrix}\right)}\SL(2,\C)$ of $\SL(2,\C)$, and compute the Lie bracket there as
\bse
[X_{\left(\begin{smallmatrix}1 & 0 \\ 0 & 1\end{smallmatrix}\right)}, Y_{\left(\begin{smallmatrix}1 & 0 \\ 0 & 1\end{smallmatrix}\right)}]_{T_{\left(\begin{smallmatrix}1 & 0 \\ 0 & 1\end{smallmatrix}\right)}\SL(2,\C)} := j^{-1} ([j(X_{\left(\begin{smallmatrix}1 & 0 \\ 0 & 1\end{smallmatrix}\right)}),j(Y_{\left(\begin{smallmatrix}1 & 0 \\ 0 & 1\end{smallmatrix}\right)})]_{\sl(2, \C)})
\ese

with 
\bse
j(X_{\left(\begin{smallmatrix}1 & 0 \\ 0 & 1\end{smallmatrix}\right)}) \left( {\left(\begin{smallmatrix}a & b \\ c & d\end{smallmatrix}\right)} \right)  = 
\Bigl( \ell_{\left(\begin{smallmatrix}a & b \\ c & d\end{smallmatrix}\right) *} \Bigr)_{\left(\begin{smallmatrix}1 & 0 \\ 0 & 1\end{smallmatrix}\right)} \Bigl( X_{\left(\begin{smallmatrix}1 & 0 \\ 0 & 1\end{smallmatrix}\right)} \Bigr)
\ese

\vspace{10pt}

So we need to compute this bracket. First thing first, for any vector $X_{\left(\begin{smallmatrix}1 & 0 \\ 0 & 1\end{smallmatrix}\right)} \in T_{\left(\begin{smallmatrix}1 & 0 \\ 0 & 1\end{smallmatrix}\right)}\SL(2,\C)$ we have to compute $j(X_{\left(\begin{smallmatrix}1 & 0 \\ 0 & 1\end{smallmatrix}\right)})$. \\

Recall that if $(U,x)$ is a chart on a manifold $M$ and $p\in U$, then the chart $(U,x)$ induces the co-ordinate induced basis which is a basis of the tangent space $T_pM$. In our case we can use the chart $(U,x)$ that contains the identity $\left(\begin{smallmatrix}1 & 0 \\ 0 & 1\end{smallmatrix}\right)$ (we must include the identity since we are interested in the tangent space at the identity $T_{\left(\begin{smallmatrix}1 & 0 \\ 0 & 1\end{smallmatrix}\right)}\SL(2,\C)$)  and hence we get an induced co-ordinate basis

\bse
\biggl\{\tvb{x}{i}{\left(\begin{smallmatrix}1 & 0 \\ 0 & 1\end{smallmatrix}\right)}\in  T_{\left(\begin{smallmatrix}1 & 0 \\ 0 & 1\end{smallmatrix}\right)}\SL(2,\C) \ \Big| \ 1\leq i \leq 3 \biggr\}
\ese

\vspace{5pt}

so that any $X_{\left(\begin{smallmatrix}1 & 0 \\ 0 & 1\end{smallmatrix}\right)} \in  T_{\left(\begin{smallmatrix}1 & 0 \\ 0 & 1\end{smallmatrix}\right)}\SL(2,\C)$ can be written as
\bse
X_{\left(\begin{smallmatrix}1 & 0 \\ 0 & 1\end{smallmatrix}\right)} =  X^i \tvb{x}{i}{\left(\begin{smallmatrix}1 & 0 \\ 0 & 1\end{smallmatrix}\right)}, \,\,\,\,\, X^i \in \C
\ese

Hence it suffices to compute the bracket only for the basis, since every other vector can be reconstructed by using the basis.\\

Let us now determine the image of these co-ordinate induced basis elements under the isomorphism $j$. The object
\bse
j\biggl(\tvb{x}{i}{\left(\begin{smallmatrix}1 & 0 \\ 0 & 1\end{smallmatrix}\right)}\biggr) \in \sl(2,\C) 
\ese

\vspace{5pt}

is a left-invariant vector field on $\SL(2,\C)$. It assigns to each point $\left(\begin{smallmatrix}a & b \\ c & d\end{smallmatrix}\right)\in U\se \SL(2,\C)$ the tangent vector
\bse
j\biggl( \tvb{x}{i}{\left(\begin{smallmatrix}1 & 0 \\ 0 & 1\end{smallmatrix}\right)}\biggr) \biggl( {\left(\begin{smallmatrix}a & b \\ c & d\end{smallmatrix}\right)} \biggr) : =
\Bigl( \ell_{\left(\begin{smallmatrix}a & b \\ c & d\end{smallmatrix}\right) *} \Bigr)_{\left(\begin{smallmatrix}1 & 0 \\ 0 & 1\end{smallmatrix}\right)} \tvb{x}{i}{\left(\begin{smallmatrix}1 & 0 \\ 0 & 1\end{smallmatrix}\right)}
\ese

\vspace{5pt}

By acting on a function $f$
\bi{rCl}
j\biggl( \tvb{x}{i}{\left(\begin{smallmatrix}1 & 0 \\ 0 & 1\end{smallmatrix}\right)}\biggr) \biggl( {\left(\begin{smallmatrix}a & b \\ c & d\end{smallmatrix}\right)} \biggr) (f) & := & \Bigl( \ell_{\left(\begin{smallmatrix}a & b \\ c & d\end{smallmatrix}\right) *} \Bigr)_{\left(\begin{smallmatrix}1 & 0 \\ 0 & 1\end{smallmatrix}\right)} \tvb{x}{i}{\left(\begin{smallmatrix}1 & 0 \\ 0 & 1\end{smallmatrix}\right)} (f)
\ei

\vspace{5pt}

by using the definition of the push-forward
\bi{rCl}
j\biggl( \tvb{x}{i}{\left(\begin{smallmatrix}1 & 0 \\ 0 & 1\end{smallmatrix}\right)}\biggr) \biggl( {\left(\begin{smallmatrix}a & b \\ c & d\end{smallmatrix}\right)} \biggr) (f) & = &\tvb{x}{i}{\left(\begin{smallmatrix}1 & 0 \\ 0 & 1\end{smallmatrix}\right)} \Bigl(f\circ \ell_{\left(\begin{smallmatrix}a & b \\ c & d\end{smallmatrix}\right)}\Bigr),
\ei

\vspace{5pt}

and by expanding the notation of $\tvb{x}{i}{\left(\begin{smallmatrix}1 & 0 \\ 0 & 1\end{smallmatrix}\right)}$ to its actual form we obtain
\bi{rCl}
j\biggl( \tvb{x}{i}{\left(\begin{smallmatrix}1 & 0 \\ 0 & 1\end{smallmatrix}\right)}\biggr) \biggl( {\left(\begin{smallmatrix}a & b \\ c & d\end{smallmatrix}\right)} \biggr) (f) & = & \partial_i\Bigl(f\circ \ell_{\left(\begin{smallmatrix}a & b \\ c & d\end{smallmatrix}\right)} \circ x^{-1}\Bigr) (x\left(\begin{smallmatrix}1 & 0 \\ 0 & 1\end{smallmatrix}\right)),
\ei

\vspace{5pt}

We can now manipulate a bit the left part of the equation by inserting an identity in the composition and we obtain
\bi{rCl}
j\biggl( \tvb{x}{i}{\left(\begin{smallmatrix}1 & 0 \\ 0 & 1\end{smallmatrix}\right)}\biggr) \biggl( {\left(\begin{smallmatrix}a & b \\ c & d\end{smallmatrix}\right)} \biggr) (f) & = &\partial_i\Bigl(f\circ {\id_U} \circ \ell_{\left(\begin{smallmatrix}a & b \\ c & d\end{smallmatrix}\right)} \circ x^{-1}\Bigr) (x\left(\begin{smallmatrix}1 & 0 \\ 0 & 1\end{smallmatrix}\right)) \\
& = & \partial_i\Bigl(f\circ ( x^{-1}\circ x) \circ \ell_{\left(\begin{smallmatrix}a & b \\ c & d\end{smallmatrix}\right)} \circ x^{-1}\Bigr) (x\left(\begin{smallmatrix}1 & 0 \\ 0 & 1\end{smallmatrix}\right))\\
& = & \partial_i\Bigl((f\circ  x^{-1})\circ (x \circ \ell_{\left(\begin{smallmatrix}a & b \\ c & d\end{smallmatrix}\right)} \circ x^{-1})\Bigr) (x\left(\begin{smallmatrix}1 & 0 \\ 0 & 1\end{smallmatrix}\right)),
\ei

\vspace{5pt}

where $f\circ  x^{-1}\cl x(U)\se \C^3 \to \C$ and $(x \circ \ell_{\left(\begin{smallmatrix}a & b \\ c & d\end{smallmatrix}\right)} \circ x^{-1})\cl x(U)\se \C^3 \to x(U)\se\C^3$ and hence, we can use the multi-dimensional chain rule to obtain
\bi{rCl}
\ldots & = & \Bigl(\partial_m(f\circ  x^{-1})\bigl((x \circ \ell_{\left(\begin{smallmatrix}a & b \\ c & d\end{smallmatrix}\right)} \circ x^{-1}) (x\left(\begin{smallmatrix}1 & 0 \\ 0 & 1\end{smallmatrix}\right))\bigr)\Bigr)\Bigl( 
\partial_i (x^m \circ \ell_{\left(\begin{smallmatrix}a & b \\ c & d\end{smallmatrix}\right)} \circ x^{-1}) (x\left(\begin{smallmatrix}1 & 0 \\ 0 & 1\end{smallmatrix}\right))\Bigr),
\ei

\vspace{5pt}

with the summation going from $m=1$ to $m=3$. \\

The first factor is simply
\bi{rCl}
(\partial_m(f\circ  x^{-1})\bigl((x \circ \ell_{\left(\begin{smallmatrix}a & b \\ c & d\end{smallmatrix}\right)} \circ x^{-1}) (x\left(\begin{smallmatrix}1 & 0 \\ 0 & 1\end{smallmatrix}\right))\bigr) & =\phantom{:} & (\partial_m(f\circ  x^{-1})\bigl((x \circ \ell_{\left(\begin{smallmatrix}a & b \\ c & d\end{smallmatrix}\right)}) x^{-1}(x\left(\begin{smallmatrix}1 & 0 \\ 0 & 1\end{smallmatrix}\right))\bigr)\\ 
& =\phantom{:} & \partial_m(f\circ  x^{-1})\bigl((x \circ \ell_{\left(\begin{smallmatrix}a & b \\ c & d\end{smallmatrix}\right)}) \left(\begin{smallmatrix}1 & 0 \\ 0 & 1\end{smallmatrix}\right)\bigr) \\
& =\phantom{:} & \partial_m(f\circ  x^{-1})(x\left(\begin{smallmatrix}a & b \\ c & d\end{smallmatrix}\right))\\
& = & \tvb{x}{m}{\left(\begin{smallmatrix}a & b \\ c & d\end{smallmatrix}\right)} (f).
\ei

\vspace{5pt}

To see what the second factor is, we first consider the map $x^m \circ \ell_{\left(\begin{smallmatrix}a & b \\ c & d\end{smallmatrix}\right)} \circ x^{-1}$. This map acts on the triple $(e,f,g)\in x(U)$ as
\bi{rCl}
(x^m \circ \ell_{\left(\begin{smallmatrix}a & b \\ c & d\end{smallmatrix}\right)} \circ x^{-1}) (e,f,g) & = & (x^m \circ \ell_{\left(\begin{smallmatrix}a & b \\ c & d\end{smallmatrix}\right)} ) \biggl(\begin{matrix}e & f \\ g & \frac{1+fg}{e}\end{matrix}\biggr)\\
& = & x^m (\biggl(\begin{matrix}a & b \\ c & d\end{matrix}\biggr) \bullet \biggl(\begin{matrix}e & f \\ g & \frac{1+fg}{e}\end{matrix}\biggr))\\
& = & x^m (\left(\begin{matrix}ae+bg &\, af+ \frac{b(1+fg)}{e} \\ ce+dg &\, cf+\frac{d(1+fg)}{e}\end{matrix}\right) ) \\
& = & (ae+bg, af+ \tfrac{b(1+fg)}{e}, ce+dg )
\ei

\vspace{5pt}

We now have to apply $\partial_i$ to this map, with $i\in \{1,2,3\}$, i.e.\ we have to differentiate with respect to each of the three complex variables $e$, $f$, and $g$. Of course depending to which derivative we take, and to which slot we apply it we end up with $3 \times 3$ different results, that we will collect in one matrix just for notation purposes. More specifically we can write the result as
\bse
\partial_i(x^m \circ \ell_{\left(\begin{smallmatrix}a & b \\ c & d\end{smallmatrix}\right)} \circ x^{-1}) (e,f,g)= D(e,f,g)^m_{\phantom{m}i},
\ese

\vspace{5pt}

where $m$ labels the rows and $i$ the columns of the matrix. (So for example the first row of the first column means the derivative of the first component (a.k.a $ae+bg$) with respect to the first variable (a.k.a $e$) hence the result (a.k.a $a$). If we do that for all components and all derivatives we end up with
\bse
D(e,f,g)= \left(\begin{matrix}a & 0 & b\\ -\frac{b(1+fg)}{e^2} &\,a+\frac{bg}{e} &\frac{bf}{e}\\ c & 0 & d\end{matrix}\right).
\ese

\vspace{5pt}

Finally, by evaluating this at $(e,f,g)=x\left(\begin{smallmatrix}1 & 0 \\ 0 & 1\end{smallmatrix}\right) = (1,0,0)$, we obtain
\bse
\partial_i(x^m \circ \ell_{\left(\begin{smallmatrix}a & b \\ c & d\end{smallmatrix}\right)} \circ x^{-1}) (x\left(\begin{smallmatrix}1 & 0 \\ 0 & 1\end{smallmatrix}\right))= D(1,0,0)^m_{\phantom{m}i},
\ese

\vspace{5pt}

where, by recalling that $d=\frac{1+bc}{a}$ we simply get,
\bse
D(1,0,0)= \left(\begin{matrix}a & 0 & b\ \\ -b & a & 0\\ c & 0 & \frac{1+bc}{a}\end{matrix}\right) := D^m_{\phantom{m}i}
\ese

\vspace{5pt}

Hence we computed the two parts of the equation above, so we can combine them and we get
\bse
j\biggl( \tvb{x}{i}{\left(\begin{smallmatrix}1 & 0 \\ 0 & 1\end{smallmatrix}\right)}\biggr) \biggl( {\left(\begin{smallmatrix}a & b \\ c & d\end{smallmatrix}\right)} \biggr) (f) =  D^m_{\phantom{m}i} \tvb{x}{m}{\left(\begin{smallmatrix}a & b \\ c & d\end{smallmatrix}\right)} (f)
\ese

\vspace{5pt}

Since this holds for an arbitrary $f\in\mathcal{C}^\infty(\SL(2,\C))$, we have
\bse
j\biggl( \tvb{x}{i}{\left(\begin{smallmatrix}1 & 0 \\ 0 & 1\end{smallmatrix}\right)}\biggr) \biggl( {\left(\begin{smallmatrix}a & b \\ c & d\end{smallmatrix}\right)} \biggr) =  D^m_{\phantom{m}i} \tvb{x}{m}{\left(\begin{smallmatrix}a & b \\ c & d\end{smallmatrix}\right)}
\ese

\vspace{5pt}

and since the point $\left(\begin{smallmatrix}a & b \\ c & d\end{smallmatrix}\right)\in U\se\SL(2,\C)$ is also arbitrary, we have
\bse
j\biggl( \tvb{x}{i}{\left(\begin{smallmatrix}1 & 0 \\ 0 & 1\end{smallmatrix}\right)}\biggr) = D^m_{\phantom{m}i}\, \frac{\partial}{\partial x^m} \in \sl(2,\C),
\ese

\vspace{5pt}

where $D$ is now the corresponding matrix of co-ordinate functions
\bse
D:=\left(\begin{matrix}x^1 & 0 & x^2\ \\ -x^2 & x^1 & 0\\ x^3 & 0 & \frac{1+x^2x^3}{x^1}\end{matrix}\right).
\ese

\vspace{5pt}

Hence, we have an expansion of the images of the basis of $T_{\left(\begin{smallmatrix}1 & 0 \\ 0 & 1\end{smallmatrix}\right)}\SL(2,\C)$ under $j$:
\bi{rCl}
j\biggl( \tvb{x}{1}{\left(\begin{smallmatrix}1 & 0 \\ 0 & 1\end{smallmatrix}\right)}\biggr) & = & x^1  \frac{\partial}{\partial x^1} - x^2\frac{\partial}{\partial x^2}  + x^3\frac{\partial}{\partial x^3} \\
j\biggl( \tvb{x}{2}{\left(\begin{smallmatrix}1 & 0 \\ 0 & 1\end{smallmatrix}\right)}\biggr) & = & x^1 \frac{\partial}{\partial x^2}  \\
j\biggl( \tvb{x}{3}{\left(\begin{smallmatrix}1 & 0 \\ 0 & 1\end{smallmatrix}\right)}\biggr) & = & x^2\, \frac{\partial}{\partial x^1} + \tfrac{1+x^2x^3}{x^1} \frac{\partial}{\partial x^3}
\ei

\vspace{5pt}

Note that while the three vector fields
\bi{rrCl}
\frac{\partial}{\partial x^m} \cl &  \SL(2,\C) & \to &  T\SL(2,\C)\\
& \biggl(\begin{matrix}a & b \\ c & d\end{matrix}\biggr) & \mapsto & \tvb{x}{m}{\left(\begin{smallmatrix}a & b \\ c & d\end{smallmatrix}\right)}
\ei
are not individually left-invariant, their linear combination with coefficients $D^m_{\phantom{m}i}$ is indeed left-invariant.\\ 

Hence we found 
\bse
j\biggl(\tvb{x}{i}{\left(\begin{smallmatrix}1 & 0 \\ 0 & 1\end{smallmatrix}\right)}\biggr) \in \sl(2,\C) 
\ese

\vspace{5pt}

We now have to calculate the bracket (in $\sl(2,\C)$) of every pair of these. We can also do them all at once, which is a good exercise in index gymnastics. We have
\bse
\biggl[j\biggl( \tvb{x}{i}{\left(\begin{smallmatrix}1 & 0 \\ 0 & 1\end{smallmatrix}\right)}\biggr) ,j\biggl( \tvb{x}{k}{\left(\begin{smallmatrix}1 & 0 \\ 0 & 1\end{smallmatrix}\right)}\biggr) \biggr] =  \left[  D^m_{\phantom{m}i}\, \frac{\partial}{\partial x^m}, D^n_{\phantom{n}k}\, \frac{\partial}{\partial x^n}\right].
\ese

\vspace{5pt}

Letting this act on an arbitrary $f\in \mathcal{C}^\infty(\SL(2,\C))$, by definition
\bse
\left[  D^m_{\phantom{m}i}\, \frac{\partial}{\partial x^m}, D^n_{\phantom{n}k}\, \frac{\partial}{\partial x^n}\right](f) :=  D^m_{\phantom{m}i}\, \frac{\partial}{\partial x^m} \Bigl( D^n_{\phantom{n}k}\, \frac{\partial}{\partial x^n} (f)\Bigr) -  D^n_{\phantom{n}k}\, \frac{\partial}{\partial x^n} \Bigl(D^m_{\phantom{m}i}\, \frac{\partial}{\partial x^m}(f)\Bigr).
\ese

\vspace{5pt}

The first term gives
\bi{rCl}
 D^m_{\phantom{m}i}\, \frac{\partial}{\partial x^m} \Bigl( D^n_{\phantom{n}k}\, \frac{\partial}{\partial x^n} (f)\Bigr) & = &  D^m_{\phantom{m}i}\, \frac{\partial}{\partial x^m} ( D^n_{\phantom{n}k}\, \partial_n (f\circ x^{-1})\circ x)\\
& = & D^m_{\phantom{m}i}\, \frac{\partial}{\partial x^m} (D^n_{\phantom{n}k})\,(\partial_n (f\circ x^{-1})\circ x) +  D^m_{\phantom{m}i}D^n_{\phantom{n}k} \,\frac{\partial}{\partial x^m} (\partial_n (f\circ x^{-1})\circ x)\\
& = & D^m_{\phantom{m}i}\, \frac{\partial}{\partial x^m} (D^n_{\phantom{n}k})\,(\partial_n (f\circ x^{-1})\circ x) +  D^m_{\phantom{m}i}D^n_{\phantom{n}k} \,\partial_m(\partial_n (f\circ x^{-1})\circ x\circ x^{-1})\circ x\\
& = & D^m_{\phantom{m}i}\, \frac{\partial}{\partial x^m} (D^n_{\phantom{n}k})\,(\partial_n (f\circ x^{-1})\circ x) +  D^m_{\phantom{m}i}D^n_{\phantom{n}k} \,\partial_m\partial_n (f\circ x^{-1})\circ x.
\ei

\vspace{5pt}

Similarly, for the second term
\bse
D^n_{\phantom{n}k}\, \frac{\partial}{\partial x^n} \Bigl(  D^m_{\phantom{m}i}\, \frac{\partial}{\partial x^m} (f)\Bigr) = D^n_{\phantom{n}k}\, \frac{\partial}{\partial x^n} (D^m_{\phantom{m}i})\,(\partial_m (f\circ x^{-1})\circ x) +  D^n_{\phantom{n}k}D^m_{\phantom{m}i}\,\partial_n\partial_m (f\circ x^{-1})\circ x.
\ese

\vspace{5pt}

Hence, recalling that $\partial_m\partial_n=\partial_n\partial_m$ by Schwarz's theorem, we have  
\bi{rCl}
\left[  D^m_{\phantom{m}i}\, \frac{\partial}{\partial x^m}, D^n_{\phantom{n}k}\, \frac{\partial}{\partial x^n}\right](f) &=&  D^m_{\phantom{m}i}\, \frac{\partial}{\partial x^m} (D^n_{\phantom{n}k})\, (\partial_n (f\circ x^{-1})\circ x) +  \cancel{D^m_{\phantom{m}i}D^n_{\phantom{n}k} \,\partial_m\partial_n (f\circ x^{-1})\circ x}\\
& & \negmedspace {} - D^n_{\phantom{n}k}\, \frac{\partial}{\partial x^n} (D^m_{\phantom{m}i})\,(\partial_m (f\circ x^{-1})\circ x) - \cancel{D^n_{\phantom{n}k}D^m_{\phantom{m}i}\,\partial_n\partial_m (f\circ x^{-1})\circ x}\\
& = & \Bigl( D^m_{\phantom{m}i}\, \frac{\partial}{\partial x^m} (D^n_{\phantom{n}k}) - D^m_{\phantom{m}k}\, \frac{\partial}{\partial x^m} (D^n_{\phantom{n}i})\Bigr)\partial_n (f\circ x^{-1})\circ x\\
& = & \Bigl( D^m_{\phantom{m}i}\, \frac{\partial}{\partial x^m} (D^n_{\phantom{n}k}) - D^m_{\phantom{m}k}\, \frac{\partial}{\partial x^m} (D^n_{\phantom{n}i})\Bigr)\frac{\partial}{\partial x^n} (f),
\ei

\vspace{5pt}

where we relabelled some dummy indices. \\

Since the $f\in\mathcal{C}^\infty(\SL(2,\C))$ was arbitrary,
\bse
\left[  D^m_{\phantom{m}i}\, \frac{\partial}{\partial x^m}, D^n_{\phantom{n}k}\, \frac{\partial}{\partial x^n}\right] =  \Bigl( D^m_{\phantom{m}i}\, \frac{\partial}{\partial x^m} (D^n_{\phantom{n}k}) - D^m_{\phantom{m}k}\, \frac{\partial}{\partial x^m} (D^n_{\phantom{n}i})\Bigr)\frac{\partial}{\partial x^n} .
\ese

\vspace{5pt}

We can now evaluate this explicitly for each basis vector. \\

For $i=1$ and $k=2$, we have
\bi{rCl}
\left[  D^m_{\phantom{m}1} \frac{\partial}{\partial x^m}, D^n_{\phantom{n}2} \frac{\partial}{\partial x^n}\right] &=&  \Bigl( \cancel{D^m_{\phantom{m}1} \frac{\partial}{\partial x^m} (D^1_{\phantom{1}2})} - D^m_{\phantom{m}2} \frac{\partial}{\partial x^m} (D^1_{\phantom{1}1})\Bigr)\frac{\partial}{\partial x^1}\\
& &\negmedspace{}+  \Bigl( D^m_{\phantom{m}1} \frac{\partial}{\partial x^m} (D^2_{\phantom{2}2}) - D^m_{\phantom{m}2} \frac{\partial}{\partial x^m} (D^2_{\phantom{2}1})\Bigr)\frac{\partial}{\partial x^2}\\
& & \negmedspace{}+ \Bigl( \cancel{D^m_{\phantom{m}1} \frac{\partial}{\partial x^m} (D^3_{\phantom{3}2})} - D^m_{\phantom{m}2} \frac{\partial}{\partial x^m} (D^3_{\phantom{3}1})\Bigr)\frac{\partial}{\partial x^3}\\
& = & -D^1_{\phantom{1}2}\frac{\partial}{\partial x^1}+(D^1_{\phantom{1}1}+D^2_{\phantom{2}2})\frac{\partial}{\partial x^2}-D^3_{\phantom{3}2}\frac{\partial}{\partial x^3}\\
& = & 2x^1 \frac{\partial}{\partial x^2}.
\ei

\vspace{5pt}

Similarly, we compute
\bi{rCl}
\left[  D^m_{\phantom{m}1} \frac{\partial}{\partial x^m}, D^n_{\phantom{n}3} \frac{\partial}{\partial x^n}\right] &=&  \Bigl( D^m_{\phantom{m}1} \frac{\partial}{\partial x^m} (D^1_{\phantom{1}3}) - D^m_{\phantom{m}3} \frac{\partial}{\partial x^m} (D^1_{\phantom{1}1})\Bigr)\frac{\partial}{\partial x^1}\\
& &\negmedspace{}+  \Bigl( \cancel{D^m_{\phantom{m}1} \frac{\partial}{\partial x^m} (D^2_{\phantom{2}3})} - D^m_{\phantom{m}3} \frac{\partial}{\partial x^m} (D^2_{\phantom{2}1})\Bigr)\frac{\partial}{\partial x^2}\\
& & \negmedspace{}+ \Bigl( D^m_{\phantom{m}1} \frac{\partial}{\partial x^m} (D^3_{\phantom{3}3}) - D^m_{\phantom{m}3} \frac{\partial}{\partial x^m} (D^3_{\phantom{3}1})\Bigr)\frac{\partial}{\partial x^3}\\
& = & -2x^2\frac{\partial}{\partial x^1}-2(\tfrac{1+x^2x^3}{x^1})\frac{\partial}{\partial x^3}
\ei

\vspace{5pt}

and
\bi{rCl}
\left[  D^m_{\phantom{m}2} \frac{\partial}{\partial x^m}, D^n_{\phantom{n}3} \frac{\partial}{\partial x^n}\right] &=&  \Bigl( D^m_{\phantom{m}2} \frac{\partial}{\partial x^m} (D^1_{\phantom{1}3}) - \cancel{D^m_{\phantom{m}3} \frac{\partial}{\partial x^m} (D^1_{\phantom{1}2})}\Bigr)\frac{\partial}{\partial x^1}\\
& &\negmedspace{}+  \Bigl( \cancel{D^m_{\phantom{m}2} \frac{\partial}{\partial x^m} (D^2_{\phantom{2}3})} - D^m_{\phantom{m}3} \frac{\partial}{\partial x^m} (D^2_{\phantom{2}2})\Bigr)\frac{\partial}{\partial x^2}\\
& & \negmedspace{}+ \Bigl( D^m_{\phantom{m}2} \frac{\partial}{\partial x^m} (D^3_{\phantom{3}3}) - \cancel{D^m_{\phantom{m}3} \frac{\partial}{\partial x^m} (D^3_{\phantom{3}2})}\Bigr)\frac{\partial}{\partial x^3}\\
& = & (D^2_{\phantom{2}1}-D^1_{\phantom{1}3})\frac{\partial}{\partial x^1}+D^2_{\phantom{2}3}\frac{\partial}{\partial x^2}-D^3_{\phantom{3}2}\frac{\partial}{\partial x^3}\\
& = & x^1 \frac{\partial}{\partial x^1}- x^2\frac{\partial}{\partial x^2} + x^3\frac{\partial}{\partial x^3},
\ei

\vspace{5pt}

where the deleted parts come from the zeros in $D^m_i$. \\

So now we know the commutator in $\sl(2, \C)$.
\bse
[j(X_{\left(\begin{smallmatrix}1 & 0 \\ 0 & 1\end{smallmatrix}\right)}),j(Y_{\left(\begin{smallmatrix}1 & 0 \\ 0 & 1\end{smallmatrix}\right)})]_{\sl(2, \C)}
\ese

\vspace{5pt}

The only thing that remains in order to fully obtain the commutator back in $T_eG$ is simply to apply $j^{-1}$, since remember that
\bse
[X_{\left(\begin{smallmatrix}1 & 0 \\ 0 & 1\end{smallmatrix}\right)}, Y_{\left(\begin{smallmatrix}1 & 0 \\ 0 & 1\end{smallmatrix}\right)}]_{T_{\left(\begin{smallmatrix}1 & 0 \\ 0 & 1\end{smallmatrix}\right)}\SL(2,\C)} := j^{-1} ([j(X_{\left(\begin{smallmatrix}1 & 0 \\ 0 & 1\end{smallmatrix}\right)}),j(Y_{\left(\begin{smallmatrix}1 & 0 \\ 0 & 1\end{smallmatrix}\right)})]_{\sl(2, \C)})
\ese

\vspace{5pt}

Hence, by applying $j^{-1}$, which is just evaluation at the identity, to these vector fields, we finally see that the induced Lie bracket on $T_{\left(\begin{smallmatrix}1 & 0 \\ 0 & 1\end{smallmatrix}\right)}\SL(2,\C)$ satisfies

\bi{rCl}
\biggl[\tvb{x}{1}{\left(\begin{smallmatrix}1 & 0 \\ 0 & 1\end{smallmatrix}\right)},\tvb{x}{2}{\left(\begin{smallmatrix}1 & 0 \\ 0 & 1\end{smallmatrix}\right)} \biggr] & = & 2\tvb{x}{2}{\left(\begin{smallmatrix}1 & 0 \\ 0 & 1\end{smallmatrix}\right)}\\[4pt]
\biggl[\tvb{x}{1}{\left(\begin{smallmatrix}1 & 0 \\ 0 & 1\end{smallmatrix}\right)},\tvb{x}{3}{\left(\begin{smallmatrix}1 & 0 \\ 0 & 1\end{smallmatrix}\right)} \biggr] & = & -2\tvb{x}{3}{\left(\begin{smallmatrix}1 & 0 \\ 0 & 1\end{smallmatrix}\right)}\\[4pt]
\biggl[\tvb{x}{2}{\left(\begin{smallmatrix}1 & 0 \\ 0 & 1\end{smallmatrix}\right)},\tvb{x}{3}{\left(\begin{smallmatrix}1 & 0 \\ 0 & 1\end{smallmatrix}\right)} \biggr] & = & \tvb{x}{1}{\left(\begin{smallmatrix}1 & 0 \\ 0 & 1\end{smallmatrix}\right)}.
\ei

\vspace{5pt}

which can be rewritten as
\bi{rCl}
\biggl[\tvb{x}{1}{\left(\begin{smallmatrix}1 & 0 \\ 0 & 1\end{smallmatrix}\right)},\tvb{x}{2}{\left(\begin{smallmatrix}1 & 0 \\ 0 & 1\end{smallmatrix}\right)} \biggr] & = & 0\tvb{x}{1}{\left(\begin{smallmatrix}1 & 0 \\ 0 & 1\end{smallmatrix}\right)} + 2\tvb{x}{2}{\left(\begin{smallmatrix}1 & 0 \\ 0 & 1\end{smallmatrix}\right)} + 0\tvb{x}{3}{\left(\begin{smallmatrix}1 & 0 \\ 0 & 1\end{smallmatrix}\right)}\\[4pt]
\biggl[\tvb{x}{1}{\left(\begin{smallmatrix}1 & 0 \\ 0 & 1\end{smallmatrix}\right)},\tvb{x}{3}{\left(\begin{smallmatrix}1 & 0 \\ 0 & 1\end{smallmatrix}\right)} \biggr] & = & 0\tvb{x}{1}{\left(\begin{smallmatrix}1 & 0 \\ 0 & 1\end{smallmatrix}\right)} + 0\tvb{x}{2}{\left(\begin{smallmatrix}1 & 0 \\ 0 & 1\end{smallmatrix}\right)} + (-2)\tvb{x}{3}{\left(\begin{smallmatrix}1 & 0 \\ 0 & 1\end{smallmatrix}\right)}\\[4pt]
\biggl[\tvb{x}{2}{\left(\begin{smallmatrix}1 & 0 \\ 0 & 1\end{smallmatrix}\right)},\tvb{x}{3}{\left(\begin{smallmatrix}1 & 0 \\ 0 & 1\end{smallmatrix}\right)} \biggr] & = & 1\tvb{x}{1}{\left(\begin{smallmatrix}1 & 0 \\ 0 & 1\end{smallmatrix}\right)} + 0\tvb{x}{2}{\left(\begin{smallmatrix}1 & 0 \\ 0 & 1\end{smallmatrix}\right)} + 0\tvb{x}{3}{\left(\begin{smallmatrix}1 & 0 \\ 0 & 1\end{smallmatrix}\right)}
\ei

\vspace{5pt}

Hence, the structure constants of $T_{\left(\begin{smallmatrix}1 & 0 \\ 0 & 1\end{smallmatrix}\right)}\SL(2,\C)$ with respect to the co-ordinate basis are
\bse
C^2_{\phantom{2}12} = 2, \qquad C^3_{\phantom{3}13}=-2,\qquad C^1_{\phantom{1}23}=1,
\ese
with all other being either zero or related to these by anti-symmetry.

However remember from the Lie algebra part of the notes that two Lie algebras $L$ and $L'$ are isomorphic if, and only if, there exists a basis of $L$ and a basis of $L'$ in which the structure constants of $L$ and $L'$ are the same.\\

Since we have already proved that $T_eG\cong_{\mathrm{Lie \, alg}}\mathcal{L}(G)$ for any Lie group $G$, we can deduce the existence of a basis $\{X^{(1)},X^{(2)},X^{(3)}\}$ of $\sl(2,\C)$ with respect to which the structure constants are those listed above .\\

\br
From now on we will be writing $\{X^{(1)},X^{(2)},X^{(3)}\}$ as $\{X_1,X_2,X_3\}$ since this is the common notation. However, do not be confused by thinking that $\{X_1,X_2,X_3\}$ are tangent vectors at points $1$, $2$ and $3$.
\er

Hence, finally we induce from the basis and the structure constants of $T_{\left(\begin{smallmatrix}1 & 0 \\ 0 & 1\end{smallmatrix}\right)}\SL(2,\C)$ a basis and structure constants of $\sl(2,\C)$ as
\bi{rCl}
[X_1,X_2] & = & 2X_2,\\
{[X_1,X_3]} & = & -2X_3,\\
{[X_2,X_3]} & = & X_1.
\ei

\br 
This is quite familiar since it coincides with our familiar commutation relations that the matrices of $\SL(2,\C)$ satisfy. However, we have to keep in mind that nothing has been mentioned regarding matrices. $\{X_1,X_2,X_3\}$ are abstract vector fields and not matrices. However, in the next chapter we will see the connection between those two when we will talk about representation theory.
\er

Now that we found the structure constants of $\sl(2,\C)$ we can use them in order to classify it. We will follow the steps we explained in the corresponding Lie algebra chapter of the notes.\\

In the basis $\{X_1,X_2,X_3\}$, the Killing form of $\sl(2,\C)$ has components
\bse
\kappa_{ij} = C^{m}_{\phantom{m}in}C^{n}_{\phantom{n}jm},
\ese
with all indices ranging from $1$ to $3$. \\

As an example, for $i=j=1$ we have
\bi{rCl}
\kappa_{11} & = & C^{m}_{\phantom{m}1n}C^{n}_{\phantom{n}1m}\\
 & = & \cancel{C^{1}_{\phantom{1}1n}C^{n}_{\phantom{n}11}} + C^{2}_{\phantom{2}1n}C^{n}_{\phantom{n}12} + C^{3}_{\phantom{3}1n}C^{n}_{\phantom{n}13}\\
 & = & C^{2}_{\phantom{2}12}C^{2}_{\phantom{2}12} + C^{3}_{\phantom{3}13}C^{3}_{\phantom{3}13}\\
 & = & 8.
\ei

Similarly we can compute all the other components. (Since $\kappa$ is symmetric, we only need to determine $\kappa_{ij}$ for $i\leq j$). \\

After computing all the components we can write them in a $3\times 3$ array, i.e
\bse
[\kappa_{ij}] = \left(\begin{matrix}8 & 0 & 0 \\ 0 & -8 & 0 \\ 0 & 0 & 8\end{matrix}\right),
\ese
which is just shorthand for 
\bse
\kappa(X_1,X_1) = 8, \qquad 
\kappa(X_2,X_2) = -8, \qquad 
\kappa(X_3,X_3) = 8, \qquad
\kappa(X_i,X_j)=0, \forall i\neq j
\ese

\bp
The Lie algebra $\sl(2,\C)$ is semi-simple.
\ep

\bq
Since the diagonal entries of $\kappa$ are all non-zero, the Killing form is non-degenerate. By Cartan's criterion, this implies that $\sl(2,\C)$ is semi-simple.
\eq

\br
There is one more thing that can be read off from the components of $\kappa$, namely, that it is an \emph{indefinite} form, i.e.\ the sign of $\kappa(X,X)$ can be positive or negative depending on which $X\in \sl(2,\C)$ we pick. \\

A result from Lie theory states that the Killing form on the Lie algebra of a compact Lie group is always negative semi-definite, i.e.\ $\kappa(X,X)$ is always negative or zero, for all $X$ in the Lie algebra. Hence, we can conclude that $\SL(2,\C)$ is not a compact Lie group.
\er

\bp
The Lie algebra $\sl(2,\C)$ is simple.
\ep

\bq
Recall that a Lie algebra is said to be simple if it contains no non-trivial ideals, and that an ideal $I$ of a Lie algebra $L$ is a Lie subalgebra of $L$ such that
\bse
\forall \, x \in I : \forall \, y\in L : \ [x,y] \in I.
\ese

Consider the ideal of $\sl(2,\C)$
\bse
I := \{\alpha X_1+\beta X_2 + \gamma X_3 \mid \alpha,\beta,\gamma \text{ restricted so that $I$ is an ideal}\}.
\ese
Since the bracket is bilinear, it suffices to check the result of bracketing an arbitrary element of $I$ with each of the basis vectors of $\sl(2,\C)$. We find
\bi{rCl}
[\alpha X_1+\beta X_2 + \gamma X_3,X_1] & = & -2\beta X_1+2\gamma X_3,\\[2pt]
{[\alpha X_1+\beta X_2 + \gamma X_3,X_2]} & = & 2 \alpha X_2 - \gamma X_1,\\[2pt]
{[\alpha X_1+\beta X_2 + \gamma X_3,X_3]} & = & -2\alpha X_3 +\beta X_1.
\ei
We need to choose $\alpha,\beta,\gamma$ so that the results always land back in $I$. Of course, we can choose $\alpha,\beta,\gamma\in\C$ and $\alpha=\beta=\gamma=0$, which correspond respectively to the trivial ideals $\sl(2,\C)$ and $0$. If none of $\alpha,\beta,\gamma$ is zero, then you can check that the right hand sides above are linearly independent, so that $I$ contains three linearly independent vectors. Since the only $n$-dimensional subspace of an $n$-dimensional vector space is the vector space itself, we have $I=L$. Thus, we are left with the following cases:
\ben[label=\roman*)]
\item if $\alpha = 0$, then $I\se \lspan_\C(\{X_2,X_3\})$ and hence we must have $\beta = \gamma = 0$ as well;
\item if $\beta = 0$, then $I\se \lspan_\C(\{X_1,X_3\})$, hence we must have $\alpha = 0$, so that in fact $I\se \lspan_\C(\{X_3\})$, and hence $\gamma = 0$ as well;
\item if $\gamma = 0$, then $I\se \lspan_\C(\{X_1,X_2\})$, hence we must have $\alpha = 0$, so that in fact $I\se \lspan_\C(\{X_2\})$, and hence $\beta = 0$ as well.
\een
In all cases, we have $I=0$. Therefore, there are no non-trivial ideals of $\sl(2,\C)$.
\eq

By observing the bracket relations of the basis elements of $\sl(2,\C)$, we can see that
\bse
H:=\lspan_\C(\{X_1\})
\ese
is a Cartan subalgebra of $\sl(2,\C)$. Indeed, for any $h\in H$, there exists a $\xi \in \C$ such that $h=\xi X_1$, and hence we have
\bi{rCl}
\ad(h) X_2 & = & \xi [X_1,X_2] = 2\xi X_2,\\
\ad(h) X_3 & = & \xi [X_1,X_3] = -2\xi X_3.
\ei
Recall that in the section on Lie algebras, we re-interpreted these eigenvalue equations in terms of functionals $\lambda_2,\lambda_3\in H^*$ 
\bi{rrClrrCl}
\lambda_2\cl & H & \xrightarrow{\sim} & \C \qquad \qquad & \lambda_3\cl & H & \xrightarrow{\sim} & \C\\
& \xi X_1 & \mapsto & 2\xi, & & \xi X_1 & \mapsto & -2\xi 
\ei
whereby
\bi{rCl}
\ad(h) X_2 & = &\lambda_2(h) X_2,\\
\ad(h) X_3 & = & \lambda_3(h) X_3.
\ei
Then, $\lambda_2$ and $\lambda_3$ are called the roots of $\sl(2,\C)$, so that the root set is $\Phi=\{\lambda_2,\lambda_3\}$. \\

Of course, we are mainly interested in a subset $\Pi\ss \Phi$ of fundamental roots, which satisfies
\ben[label=\roman*)]
\item $\Pi$ is a linearly independent subset of $H^*$;
\item for any $\lambda\in \Phi$, we have $\lambda \in \lspan_{\epsilon,\N}(\Pi)$.
\een
We can choose $\Pi:=\{\lambda_2\}$, even though $\Pi:=\{\lambda_3\}$ would work just as well.  \\

Since $|\Pi|=1$, the Weyl group is generated by the single Weyl transformation
\bi{rrCl}
s_{\lambda_2}\cl & H^*_\R & \to & H^*_\R\\
& \mu & \mapsto & \mu - 2\frac{\kappa^*(\lambda_2,\mu)}{\kappa^*(\lambda_2,\lambda_2)}\lambda_2 .
\ei
Recall that we can recover the entire root set $\Phi$ by acting on the fundamental roots with Weyl transformations. Indeed, we have
\bse
s_{\lambda_2}(\lambda_2) = \lambda_2- 2\frac{\kappa^*(\lambda_2,\lambda_2)}{\kappa^*(\lambda_2,\lambda_2)}\lambda_2 = \lambda_2-2\lambda_2 = -\lambda_2 = \lambda_3,
\ese
as expected. \\

Since there is only one fundamental root, the Cartan matrix is actually just a $1\times 1$ matrix. Its only entry is a diagonal entry, and since $\sl(2,\C)$ is simple,
 we have
\bse
C = (2). 
\ese
The Dynkin diagram of $\sl(2,\C)$ is simply 
\begin{center}
\begin{tikzpicture}
\draw[fill=white] (0,0) circle[radius=0.15];
\end{tikzpicture}
\end{center}
Hence, with reference to the Cartan classification, we have $ \sl(2,\C) = A_1$.